\documentclass[11pt]{article}
\usepackage{amsmath,amssymb,amsthm,hyperref}
\newtheorem{theorem}{Theorem}
\newtheorem{corollary}[theorem]{Corollary}
\title{Permutation Sheet Readers: Two Exact Capacity Obstructions}
\author{}
\date{}
\begin{document}
\maketitle

\paragraph{Status.}
These are complete written auxiliary results from the resumed global
construction search. They are not Lean declarations and prove no new
asymptotic upper or lower bound for unrestricted head complexity.
The original model and denominator conditions are unchanged.
Their purpose is to close two specific attempted decoders for the regular
permutation sheets constructed in \texttt{compact-cauchy-sheets.tex}.
Factorially many regular parameter solutions do not imply factorially
many independently readable sign patterns.

\begin{theorem}[A bounded row-pole library is a linear feature map]
Let $\lambda_1,\ldots,\lambda_H$ be distinct real numbers. For each
slot $h$, fix a finite set $D_h$ of real numbers strictly larger than
$\max_j\lambda_j$. Suppose every query on a permutation $\pi$ is of
the form
\[
 Q_i(\pi)=c_i+\sum_{h=1}^H
 \frac{\alpha_{ih}+\beta_{ih}\lambda_{\pi(h)}}
      {d_{ih}-\lambda_{\pi(h)}},\qquad d_{ih}\in D_h.
\]
If $r$ fixed queries realize all $2^r$ strict sign vectors as $\pi$
varies, then $r\le\sum_h|D_h|$.
If all $D_h$ are the same singleton, then $r\le H-1$.
\end{theorem}

\begin{proof}
The identity
\[
 \frac{\alpha+\beta t}{d-t}
 =-\beta+\frac{\alpha+\beta d}{d-t}
\]
makes every query an affine functional of the
$q=\sum_h|D_h|$ coordinates
\[
 z_{h,d}(\pi)=\frac1{d-\lambda_{\pi(h)}}.
\]
The image of all permutations is a subset of $\mathbb R^q$.
An arrangement of $r$ affine hyperplanes in affine dimension $q$
has at most $\sum_{j=0}^{q}\binom rj$ open regions. This follows by
induction on $r$: inserting the last hyperplane creates at most as
many additional regions as the preceding hyperplanes induce on that
hyperplane, whose dimension is at most $q-1$. The initial cases are
one region for $r=0$ and at most one region for $q=0$.
Consequently $r>q$ leaves strictly fewer than $2^r$ strict sign
vectors even in the whole ambient space.

For one common pole $d$, the feature points satisfy the additional
affine identity
\[
 \sum_{h=1}^H z_{h,d}(\pi)=\sum_{j=1}^H\frac1{d-\lambda_j}.
\]
Their affine dimension is therefore at most $H-1$, and the same
argument gives the sharper assertion.
\end{proof}

\begin{theorem}[Subset Stieltjes readers have linear capacity]
Fix distinct real numbers $\lambda_1,\ldots,\lambda_L$, an integer
$1\le k\le L$, and a family $\mathcal S$ of $k$-element subsets of
$\{1,\ldots,L\}$. For $S\in\mathcal S$ put
\[
 F_S(d)=\sum_{j\in S}\frac1{d-\lambda_j},
 \qquad d>\max_j\lambda_j.
\]
For fixed probe pairs $(d_i,\theta_i)$ with $d_i>\max_j\lambda_j$,
if the signs of $F_S(d_i)-\theta_i$, as $S$ varies, realize all
$2^r$ strict patterns, then $r\le2k-1$.
\end{theorem}

\begin{proof}
First, repeated probe locations cannot occur in a shattered list.
Indeed, at one fixed $d$, two unequal thresholds have one impossible
strict labeling, and two equal thresholds have both mixed labelings
impossible. Thus order the locations as $d_1<\cdots<d_r$.

Suppose $r\ge2k$. Choose witnesses $S,T$ for the two complementary
alternating labelings. At every probe, $F_S(d_i)-F_T(d_i)$ has the
prescribed alternating strict sign, independently of the values of
the thresholds. In particular $S\ne T$.

Let $A=S\setminus T$, $B=T\setminus S$, and $q=|A|=|B|$, so
$1\le q\le k$. Then
\[
 F_S(d)-F_T(d)=\sum_{j\in A}\frac1{d-\lambda_j}
                 -\sum_{j\in B}\frac1{d-\lambda_j}.
\]
Multiplying by the positive polynomial
$R(d)=\prod_{j\in A\cup B}(d-\lambda_j)$ gives a polynomial $P$.
The potential coefficient of $d^{2q-1}$ is $q-q=0$, so
$\deg P\le2q-2$. The polynomial is nonzero: at $d=\lambda_a$ for
any $a\in A$, its value is
$\prod_{j\in(A\cup B)\setminus\{a\}}(\lambda_a-\lambda_j)\ne0$.
The alternating strict signs at the ordered probes, together with
continuity, give at least $r-1\ge2k-1>2q-2$ distinct real roots.
This contradicts the polynomial root bound.
\end{proof}

\begin{corollary}[Independent block moment decoding does not amplify]
Partition $H$ permutation slots into blocks of sizes $k_1,\ldots,k_b$.
Suppose each query probes the unweighted Stieltjes sum of the values
assigned to one specified block, with its own pole and threshold.
Then a shattered query family has at most
$\sum_{a=1}^b(2k_a-1)=2H-b$ members.
\end{corollary}

\begin{proof}
Every subfamily associated with one block must itself be shattered.
Its possible values are some family of $k_a$-subsets of the fixed
permutation values, so the preceding theorem applies to it. Sum the
resulting bounds over the blocks. Requiring the blocks to be disjoint
can only restrict the possible simultaneous patterns; the argument
does not presume their independent realizability.
\end{proof}

\paragraph{Exact transfer boundary.}
The second theorem permits the poles to vary freely between probes
and permits any fixed threshold at each probe. It therefore also
covers proposals which first constrain low power sums of each block
and then use higher moments of its resolvent to read extra bits:
restricting the allowed subsets does not evade the bound. It does
not cover arbitrary slot-dependent weights or independently varying
row-and-slot poles. Those freedoms are the remaining multivariate
decoder question, not a consequence of the sheet count.

These obstructions concern readers on an already admissible compact
parameter family. They do not claim that an arbitrary $H$-head score
has either reader form, and hence do not imply an unrestricted lower
bound for $W(n)$.
\end{document}
